% part: intuitionistic-logic
% chapter: soundness-completeness
% section: soundness-nd

\documentclass[../../../include/open-logic-section]{subfiles}

\begin{document}

\olfileid{int}{sc}{snd}

\olsection{د طبيعي استنتاج سموالے}

اوس به د اړيکيزې معنايي پوهنې له مخې د طبيعي استنتاج سموالے ثابت
کړو: يعنې وښيو چې که !!a{formula} د فرضونو له يوه سټ څخه
!!{derivable} وي، نو د فرضونو سټ هغه !!{formula} لازموي۔

\begin{thm}[سموالے]
  \ollabel{thm:soundness}
  که $\Gamma \Proves !A$ وي، نو $ \Gamma \Entails !A$ وي۔
\end{thm}

\begin{proof}
  ثابتوو چې که $\Gamma \Proves !A$ وي، نو $\Gamma \Entails !A$
  وي۔ ثبوت له~$\Gamma$ څخه د~$!A$ پر !!{derivation} استقرا ده۔
  
  \begin{enumerate}
  \item که !!{derivation} يوازې له فرض~$!A$ څخه جوړ وي،
    $!A \Proves !A$ لرو او بايد $!A \Entails !A$ وښيو۔ فرض کړئ
    $\mSat{M}{!A}[w]$ وي۔ نو په څرګنده $\mSat{M}{!A}[w]$ وي۔

  \item !!{derivation} په $\Intro{\land}$ پاے ته رسېږي:
    \iftag{probAnd}{تمرين۔}{له !!{undischarged} فرضونو~$\Gamma$
      څخه د مقدمې~$!B$ او له !!{undischarged} فرضونو~$\Delta$
      څخه د مقدمې~$!C$ !!{derivation}s ښيي چې $\Gamma \Proves
      !B$ او $\Delta \Proves !C$۔ د استقرا د فرض له مخې
      $\Gamma \Entails !B$ او $\Delta \Entails !C$ لرو۔ بايد
      $\Gamma \cup \Delta \Entails !B \land !C$ وښيو، ځکه د ټول
      !!{derivation} !!{undischarged} فرضونه د $\Gamma$ او
      $\Delta$ ګډ سټ دے۔ نو $\mSat{M}{\Gamma
        \cup \Delta}[w]$ فرض کړئ۔ بيا $\mSat{M}{\Gamma}[w]$
      هم وي۔ ځکه چې $\Gamma \Entails !B$، نو $\mSat{M}{!B}[w]$
      وي۔ په ورته ډول $\mSat{M}{!C}[w]$ وي۔ نو
      $\mSat{M}{!B \land !C}[w]$ وي۔ د ژباړې د نسخې څرګندونه
      (OLINT-007): منجمده سرچينه په دې حالت کښې د مقدمې د دوو
      بېلو نښو پر ځاے د بل عطف پايله ليکي؛ د همدغو مقدمو او
      وروستۍ کرښې له مخې دلته د هغو عطف مراد دے۔}

  \item !!{derivation} په $\Elim{\land}$ پاے ته رسېږي:
    \iftag{probAnd}{تمرين۔}{له !!{undischarged} فرضونو $\Gamma$
        څخه د مقدمې $!B \land !C$ !!{derivation} ښيي چې
        $\Gamma \Proves !B \land !C$۔ د استقرا د فرض له مخې
        $\Gamma \Entails !B \land !C$۔ بايد $\Gamma
        \Entails !B$ وښيو۔ نو $\mSat{M}{\Gamma}[w]$ فرض کړئ۔
        ځکه $\Gamma \Entails !B \land !C$، نو
        $\mSat{M}{!B \land !C}[w]$ وي۔ بيا
        $\mSat{M}{!B}[w]$ هم وي۔ په ورته ډول که $\Elim\land$
        په~$!C$ پاے ته ورسېږي، نو $\Gamma \Entails !C$ وي۔}

  \item !!{derivation} په $\Intro{\lor}$ پاے ته رسېږي:
    \iftag{probOr}{تمرين۔}{فرض کړئ مقدمه $!B$ ده او په~$!B$
      ختمېدونکي !!{derivation} !!{undischarged} فرضونه $\Gamma$
      دي۔ نو $\Gamma \Proves !B$ لرو او د استقرا د فرض له مخې
      $\Gamma \Entails !B$۔ بايد $\Gamma \Entails !B \lor !C$
      وښيو۔ $\mSat{M}{\Gamma}[w]$ فرض کړئ۔ ځکه
      $\Gamma \Entails !B$، نو $\mSat{M}{!B}[w]$ وي۔ بيا
      $\mSat{M}{!B \lor !C}[w]$ هم وي۔ په ورته ډول، که مقدمه
      ~$!C$ وي، نو $\Gamma \Entails !C$ لرو۔}
    
  \item !!{derivation} په~$\Elim{\lor}$ پاے ته رسېږي:
    \iftag{probOr}{تمرين۔}{په مقدمو ختمېدونکي !!{derivation}s دا دي:
      $!B \lor !C$ له !!{undischarged} فرضونو~$\Gamma$ څخه،
      $!D$ له !!{undischarged} فرضونو
      $\Delta_1 \cup \{!B\}$ څخه، او $!D$ له !!{undischarged}
      فرضونو~$\Delta_2 \cup \{!C\}$ څخه۔ نو $\Gamma \Proves
      !B \lor !C$، $\Delta_1 \cup \{!B\} \Proves !D$ او
      $\Delta_2 \cup \{!C\} \Proves !D$ لرو۔ د استقرا د فرض له
      مخې $\Gamma \Entails !B \lor !C$،
      $\Delta_1 \cup \{!B\} \Entails !D$ او
      $\Delta_2 \cup \{!C\} \Entails !D$ دي۔ بايد
      $\Gamma \cup \Delta_1 \cup \Delta_2 \Entails !D$ ثابت کړو۔

    $\mSat{M}{\Gamma \cup \Delta_1 \cup \Delta_2}[w]$ فرض کړئ۔
    نو $\mSat{M}{\Gamma}[w]$ وي، او ځکه چې
    $\Gamma \Entails !B \lor !C$، نو
    $\mSat{M}{!B \lor !C}[w]$ وي۔ د $\mSat{M}{}$ د تعريف له
    مخې يا $\mSat{M}{!B}[w]$ وي يا $\mSat{M}{!C}[w]$۔ نو دوه
    حالتونه بېلوو: (الف)~$\mSat{M}{!B}[w]$ وي۔ بيا
    $\mSat{M}{\Delta_1 \cup \{!B\}}[w]$ وي۔ ځکه
    $\Delta_1 \cup \{!B\} \Entails !D$، نو
    $\mSat{M}{!D}[w]$ وي۔ (ب)~$\mSat{M}{!C}[w]$ وي۔ بيا
    $\mSat{M}{\Delta_2 \cup \{!C\}}[w]$ وي۔ ځکه
    $\Delta_2 \cup \{!C\} \Entails !D$، نو
    $\mSat{M}{!D}[w]$ وي۔ په دواړو حالاتو کښې
    $\mSat{M}{!D}[w]$ لرو، لکه چې غوښتل مو۔
    د ژباړې د نسخې څرګندونه (OLINT-008): منجمده سرچينه د لومړي
    حالت د صدق نښه ناسمه تړي، او په دواړو حالتونو کښې د فرض
    يوازېنی غړے له سټه پرته ليکي؛ دلته د مخکينيو فرضونو له
    سټونو سره د دواړو حالتونو نښې برابرې شوې دي۔}

  \item !!{derivation} په $\Intro{\lif}$ پاے ته رسېږي او
    پايله ئې~$!B\lif !C$ ده۔ نو مقدمه~$!C$ ده، او په دې مقدمه
    ختمېدونکے !!{derivation} !!{undischarged} فرضونه~$\Gamma \cup
    \{!B\}$ لري۔ په دې توګه $\Gamma \cup \{!B\} \Proves !C$،
    او د استقرا د فرض له مخې $\Gamma \cup \{!B\} \Entails !C$
    لرو۔ بايد $\Gamma \Entails !B \lif !C$ وښيو۔

    $\mSat{M}{\Gamma}[w]$ فرض کړئ۔ بايد وښيو چې د هر داسې
    $w'$ لپاره چې $Rww'$ وي، که $\mSat{M}{!B}[w']$ وي، نو
    $\mSat{M}{!C}[w']$ وي۔ نو $Rww'$ او
    $\mSat{M}{!B}[w']$ فرض کړئ۔ د
    \olref[sem][rel]{prop:true-monotonic} له مخې
    $\mSat{M}{\Gamma}[w']$ وي۔ ځکه
    $\Gamma \cup \{!B\} \Entails !C$، نو
    $\mSat{M}{!C}[w']$ وي؛ همدا مو غوښتل۔

  \item !!{derivation} په $\Elim{\lif}$ پاے ته رسېږي او
    پايله ئې~$!C$ ده۔ مقدمې $!B \lif !C$ او $!B$ دي چې په
    ترتيب له !!{undischarged} فرضونو $\Gamma$ او $\Delta$
    څخه !!{derivation}s لري۔ نو $\Gamma \Proves !B \lif !C$ او
    $\Delta \Proves !B$ لرو۔ د استقرا د فرض له مخې
    $\Gamma \Entails !B \lif !C$ او $\Delta \Entails !B$
    دي۔ بايد $\Gamma \cup \Delta \Entails !C$ وښيو۔

    $\mSat{M}{\Gamma \cup \Delta}[w]$ فرض کړئ۔ ځکه
    $\mSat{M}{\Gamma}[w]$ او $\Gamma \Entails !B \lif !C$
    دي، نو $\mSat{M}{!B \lif !C}[w]$ وي۔ د تعريف له مخې د هر
    داسې $w'$ لپاره چې $Rww'$ وي، که $\mSat{M}{!B}[w']$ وي،
    نو $\mSat{M}{!C}[w']$ وي۔ څرنګه چې $R$ انعکاسي ده، $w$
    هم په هغو $w'$ کښې دے چې $Rww'$ لري؛ يعنې که
    $\mSat{M}{!B}[w]$ وي، نو $\mSat{M}{!C}[w]$ وي۔ ځکه
    $\mSat{M}{\Delta}[w]$ او $\Delta \Entails !B$ دي، نو
    $\mSat{M}{!B}[w]$ وي۔ نو $\mSat{M}{!C}[w]$ شو؛ همدا مو
    غوښتل۔

  \item !!{derivation} په $\FalseInt$ پاے ته رسېږي او پايله ئې
    ~$!A$ ده۔ مقدمه $\lfalse$ ده او د مقدمې د !!{derivation}
    !!{undischarged} فرضونه~$\Gamma$ دي۔ نو $\Gamma \Proves
    \lfalse$ وي۔ د استقرا د فرض له مخې
    $\Gamma \Entails \lfalse$ وي۔ بايد $\Gamma \Entails !A$
    وښيو۔

    په نامستقيم ډول ځو۔ که $\Gamma \Entails/ !A$ وي، داسې
    مدل~$\mModel{M}$ او نړۍ~$w$ شته چې
    $\mSat{M}{\Gamma}[w]$ او $\mSat/{M}{!A}[w]$ وي۔ ځکه
    $\Gamma \Entails \lfalse$، نو $\mSat{M}{\lfalse}[w]$
    وي۔ خو دا ناشونې ده، ځکه د تعريف له مخې
    $\mSat/{M}{\lfalse}[w]$ وي۔ نو $\Gamma \Entails !A$ دے۔
  \item !!{derivation} په $\Intro\lnot$ پاے ته رسېږي: تمرين۔
  \item !!{derivation} په $\Elim\lnot$ پاے ته رسېږي: تمرين۔
  \end{enumerate}
\end{proof}

\begin{prob}
  د \olref[int][sc][snd]{thm:soundness} ثبوت بشپړ کړئ۔ د
  $\Intro\lnot$ او $\Elim\lnot$ د حالتونو لپاره د
  \olref[int][sem][rel]{defn:true-at-w} له مخې د
  $\mSat{M}{\lnot !A}[w]$ تعريف وکاروئ؛ يعنې $\lnot !A$ د
  $!A \lif \lfalse$ په توګه تعريف شوے مۀ ګڼئ۔
\end{prob}

\begin{prob}
  وښيئ چې لاندې !!{formula}s په شهودي منطق کښې !!{derivable} نۀ دي:
  \begin{enumerate}
    \item $(!A \lif !B) \lor (!B \lif !A)$
    \item $(\lnot\lnot !A \lif !A) \lif (!A \lor \lnot !A)$
    \item $(!A \lif !B \lor !C) \lif \bigl((!A \lif !B)\lor(!A \lif !C)\bigr)$
  \end{enumerate}
\end{prob}

\end{document}
